\documentclass[journal,compsoc]{IEEEtran}

\usepackage{amsmath,amssymb,amsfonts}
\usepackage{algorithmic}
\usepackage{graphicx}
\usepackage{textcomp}
\usepackage{xcolor}
\usepackage{booktabs}
\usepackage{url}
\usepackage{hyperref}
\usepackage{cite}

\hypersetup{
    colorlinks=true,
    linkcolor=blue,
    filecolor=magenta,      
    urlcolor=cyan,
    citecolor=red,
}

\begin{document}

\title{Continuous Authorization over the Agentic-RAG Context Lifecycle}

\author{Taha~Bin~Hanif%
\thanks{Taha Bin Hanif is with the PNTLab, CNIT, and the TeCIP Institute, Scuola Superiore Sant'Anna, Pisa, Italy (e-mail: taha.binhanif@santannapisa.it).}}

\IEEEtitleabstractindextext{%
\begin{abstract}
Enterprise retrieval-augmented generation (RAG) and autonomous agent frameworks increasingly process sensitive multi-tenant documents and dispatch high-consequence tools. While recent authorization-first retrieval (AFR) designs prevent cross-tenant exposure during initial index queries, they remain strictly point-in-time mechanisms. We demonstrate that AFR-only architectures fail under lifecycle threats: when access policies are revoked across multi-turn sessions, stateless conversation memory causes 100\% stale unauthorized context exposure. Furthermore, indirect prompt injections embedded in authorized documents can hijack model-selected tool calls. In this paper, we propose a multi-point continuous authorization architecture that enforces fine-grained Relationship-Based Access Control (ReBAC via SpiceDB) across three trust boundaries: pre-retrieval vector index filtering, retained-context revalidation before every cross-turn reuse, and action-time tool dispatch. To evaluate this approach, we release the \textsc{AuthInject-Lifecycle} benchmark, comprising 160 multi-turn test instances across seven attack families. In factorial evaluations across three frontier LLMs (LLaMA-3.3-70B, DeepSeek-R1-32B, GPT-4o-Mini) and twelve configurations, our continuous authorization architecture achieves 0.0\% unauthorized and structural context exposure, eliminates stale ACL bypasses (compared to 100\% in AFR-only), and reduces model-chosen tool hijacking from 46.2\% to 4.4\%, at an acceptable +42\% median latency overhead.
\end{abstract}

\begin{IEEEkeywords}
Agentic RAG, Continuous Authorization, Relationship-Based Access Control, Indirect Prompt Injection, Multi-Tenant Security, Tool-Use Security.
\end{IEEEkeywords}}

\maketitle
\IEEEdisplaynontitleabstractindextext
\IEEEpeerreviewmaketitle

\section{Introduction}
\label{sec:intro}

Retrieval-Augmented Generation (RAG) grounds Large Language Models (LLMs) in external knowledge bases, mitigating factual hallucinations. In enterprise deployments, however, RAG pipelines operate over confidential multi-tenant document stores and integrate autonomous agentic loops capable of invoking external APIs and database tools. This convergence creates a critical security vulnerability: standard semantic retrieval conflates relevance with authorization, and autonomous agent loops execute untrusted retrieved instructions as operational mandates.

\subsection{Background and the Point-in-Time Limitations of AFR}
\label{sec:intro:background}

To prevent multi-tenant data leakage, post-retrieval filtering (retrieve-then-filter) was historically employed. However, post-filtering remains susceptible to semantic bait attacks, where unauthorized documents displace legitimate context in the similarity ranking. To address this, Namboothiri et al.~\cite{namboothiri2024author} formalized Authorization-First Retrieval (AFR) candidate restriction as a retrieval-time noninterference invariant, pushing access control filters into the vector database index.

While AFR successfully secures the initial retrieval step, enterprise agentic RAG operates over extended multi-turn dialogs where context is cached, transformed, and reused across turns. Our key insight is that \textit{authorization is a continuous lifecycle invariant, not a point-in-time retrieval gate}. When a user's document permissions are revoked mid-session (e.g., role change or document declassification), standard agent conversation buffers retain unauthorized chunks, enabling trivial stale-context leakage. Furthermore, when authorized documents contain indirect prompt injections, AFR provides zero defense against malicious tool invocation.

\subsection{Research Questions}
\label{sec:intro:rqs}

We guide our investigation through three core research questions:
\begin{itemize}
    \item \textbf{RQ1 (Lifecycle Defense Efficacy):} Does continuous authorization across retrieval, memory, and action reduce composite failure rates relative to AFR-only under multi-turn lifecycle threats?
    \item \textbf{RQ2 (Ablation and Overhead Tradeoffs):} Which ablated security components contribute most significantly to defense-in-depth and the protection-latency tradeoff?
    \item \textbf{RQ3 (Cross-Model Generalization):} Are continuous authorization guarantees robust and statistically consistent across diverse frontier model families?
\end{itemize}

\subsection{Contributions}
\label{sec:intro:contributions}

In this paper we make the following four contributions:
\begin{enumerate}
    \item \textbf{C1 --- Multi-Point Reference Monitor Architecture:} The first system design to enforce authorization at three distinct trust boundaries: (i) pre-retrieval candidate restriction (via SpiceDB ReBAC + Qdrant \texttt{MatchAny} operationalizing~\cite{namboothiri2024author}), (ii) retained-context revalidation before every cross-turn reuse, and (iii) action-time tool-dispatch authorization. We formalize the combined invariant as \textit{continuous authorization}.
    \item \textbf{C2 --- \textsc{AuthInject-Lifecycle} Benchmark:} A 160-case, 7-threat-family combined corpus extending InjecAgent, AgentDojo, OGX, and BIPIA benchmarks with two new probe families: same-tenant disjoint-document semantic bait and stale-ACL cross-turn revocation. The benchmark enables direct comparison of post-filter, AFR-only, and continuous-authorization systems on a shared ground truth.
    \item \textbf{C3 --- Factorial Empirical Evaluation:} A 12-configuration empirical study across 4 baselines, 2 proposed systems, and 6 ablations evaluated on 3 LLMs with 5 stochastic repeats, using Wilson 95\% confidence intervals and McNemar paired significance testing. On 40 stale-ACL cases we show the prior SOTA (AFR-only~\cite{namboothiri2024author}) records 100\% structural exposure while our proposed continuous system records 0.0\% ($p = 1.28 \times 10^{-8}$).
    \item \textbf{C4 --- Open-Source Production Package:} An enterprise-grade, reproducible implementation (\texttt{secure-rag==0.3.0}) featuring Docker Compose, Kubernetes StatefulSet, and Helm deployment charts with dual-sink HMAC-SHA256 tamper-evident audit logging and automated CI validation.
\end{enumerate}

\section{Related Work}
\label{sec:related_work}

Securing autonomous, multi-tenant agentic Retrieval-Augmented Generation (RAG) systems requires unifying four historically distinct domains: enterprise authorization engines, adversarial prompt injection defenses, capability-based tool access control, and stateful multi-turn memory revocation. We review prior work across these four axes and highlight the architectural gaps addressed by our continuous authorization framework.

\subsection{Authorization and Access Control in Multi-Tenant RAG}
\label{sec:rel_rag_authz}

Enterprise multi-tenant information systems increasingly adopt Relationship-Based Access Control (ReBAC) architectures inspired by Google's Zanzibar~\cite{pang2019zanzibar}. Open-source operationalizations such as SpiceDB~\cite{authzed2023spicedb} and Open Policy Agent (OPA)~\cite{styra2022opa} provide expressive graph-traversal semantics to evaluate fine-grained permissions (e.g., direct ownership, delegated viewership, and tenant-level inheritance) with sub-millisecond query evaluation latency~\cite{rebacenterprise2024, authzbench2025}. However, classical access control models assume deterministic database querying where security filters can be compiled into structured SQL predicates.

When adapted to vector search engines~\cite{qdrant2023, lewis2020rag}, early enterprise RAG deployments implemented \textit{post-filtering} architectures~\cite{jeong2025permission}, wherein vector indexes retrieve the top-$k$ nearest semantic neighbors across the entire document corpus, after which an external authorization engine strips out unauthorized documents before presenting the context to the LLM. While computationally simple, post-filtering fundamentally violates classical noninterference principles~\cite{rushby1982noninterference, saltzer1975protection}: unauthorized document metadata and chunk embeddings remain active candidate vectors in the vector search space, introducing severe structural information exposure and cache side-channels~\cite{multitenantrag2025, authzrag2026}.

To resolve this structural vulnerability, Namboothiri et al.~\cite{namboothiri2026afr} introduced \textit{Authorization-First Retrieval} (AFR), formalizing the noninterference invariant wherein the vector retrieval engine restricts the search space \textit{a priori} to the exact set of authorized document identifiers returned by an initial ReBAC lookup. This ensures that unauthorized chunks are structurally invisible during dense passage retrieval~\cite{karpukhin2020dpr}.

Concurrently, researchers have investigated the interaction between retrieval boundaries and document poisoning. SafeRAG~\cite{liang2025saferag} and PoisonedRAG~\cite{zou2025poisonedrag} benchmarked vulnerabilities where malicious actors poison shared enterprise corpora to hijack model outputs. PrivRAG~\cite{privrag2025acl} and GuardRAG~\cite{guardrag2025} developed differential privacy and safety filtering heuristics to mitigate data extraction, while Arceo and Narsing~\cite{arceo2026securing} demonstrated the persistent ``relevance-authorization gap'' in agentic information systems when document permissions change dynamically.

\noindent\textbf{Gap in Prior Work:} Despite these advances, all prior AFR frameworks protect only the single-turn retrieval boundary. They fail to track the authorization provenance of chunks across agent reasoning loops, leaving systems vulnerable when permissions are revoked between conversational turns or when retrieved text coerces the agent into unauthorized tool execution. Our work bridges this gap by extending AFR into a continuous, end-to-end lifecycle reference monitor.

\subsection{Indirect Prompt Injection Benchmarks and Defenses}
\label{sec:rel_injection}

Unlike direct jailbreaking, where an adversarial user directly instructs the model to violate its safety alignment, \textit{Indirect Prompt Injection} (IPI)~\cite{greshake2023promptinject, sha2024promptinjection} embeds adversarial instructions within untrusted external data retrieved during task execution (e.g., search results, web pages, or tenant documents). When ingested into the LLM's context window, these passive instructions hijack the model's control flow, compelling it to leak confidential data or execute unapproved actions.

Systematic benchmarking of IPI vulnerabilities has evolved rapidly. Zhan et al.~\cite{zhan2024injecagent} introduced InjecAgent, formalizing 1,054 test cases categorizing direct-harm versus data-stealing attacks against tool-integrated models. Debenedetti et al.~\cite{debenedetti2024agentdojo} developed AgentDojo, providing a dynamic execution environment measuring environment-state attack success rates (ASR). Concurrently, AgentSecBench~\cite{zhang2025agentsecbench} and BIPIA~\cite{yi2025bipia} established multi-vector benchmark suites across diverse application domains. Recent work on adaptive injection, such as TopicAttack~\cite{topicattack2025}, has shown that heuristic keyword filters are easily evaded by semantic paraphrasing.

Defenses against IPI generally fall into three categories: (i) input classifiers and guardrail models (e.g., NeMo Guardrails~\cite{rebedea2023nemo}, PromptGuard~\cite{promptguard2024}, and LLM Guard~\cite{llmguard2024}), which attempt to classify text chunks prior to generation; (ii) test-time perturbation and alignment techniques, such as IPIGuard~\cite{ipiguard2025} and SecAlign~\cite{secalign2025ccs}; and (iii) structural context isolation, datamarking, and XML delimiters~\cite{melon2025icml, drift2025, spring2025acl}.

\noindent\textbf{Gap in Prior Work:} Existing IPI defenses treat retrieved context as static, stateless text. They do not maintain cryptographic or provenance-bound links between the retrieved chunk, its access policy version, and the executing principal. Consequently, existing IPI defenses cannot detect when an authorized chunk from turn $T_0$ becomes an unauthorized injection vector at turn $T_1$ after an ACL revocation event.

\subsection{Agent Tool-Action Authorization and Defenses}
\label{sec:rel_tools}

Autonomous agents leveraging frameworks like ReAct~\cite{yao2023react} and LangGraph~\cite{langgraph2024} autonomously plan multi-step workflows, invoke external APIs, and mutate environment state. However, granting language models tool-execution capabilities creates critical confused-deputy vulnerabilities when the agent processes untrusted retrieved documents~\cite{agentpoison2024, toolattack2025}.

Recent literature has proposed various runtime guardrails for agent actions. Task Shield~\cite{taskshield2025acl} introduced task-boundary enforcement to prevent multi-turn tool-calling deviations. TOOLSHIELD~\cite{toolshield2025} and ToolSec~\cite{toolsec2025} investigated hardware-enforced and sandbox-level capability checks for LLM tools. Furthermore, Google DeepMind's safety framework~\cite{deepmindagents2025} and Kumar et al.~\cite{toolsafety2025} formalized access control principles for foundation model tool execution.

\noindent\textbf{Gap in Prior Work:} Existing tool-defense mechanisms evaluate permissions solely at tool-dispatch time based on the agent's current token stream. They lack integration with the data-retrieval layer and therefore cannot determine whether the prompt that induced the tool action originated from authorized context or from an adversarial injection bait. Our architecture enforces action-time ReBAC reference monitoring before any side-effect occurs, strictly validating the caller's principal identity against Zanzibar policy tuples.

\subsection{Context Persistence Across Turns and Memory Revocation}
\label{sec:rel_memory}

Stateful multi-turn conversational agents frequently persist conversational history, entity summaries, and prior retrieved contexts across interaction turns~\cite{memoryrag2025}. In enterprise environments, document permissions are dynamic: employees change roles, project memberships expire, and security holds are lifted~\cite{dbsync2024vldb, nist80053ac}.

When an authorized user retrieves a document at time $T_0$ and the document's access policy is subsequently revoked at time $T_1$, naive agent implementations continue to ground subsequent answers on the cached context carried over from turn $T_0$. While distributed databases have studied cache invalidation and immediate revocation semantics extensively~\cite{dbsync2024vldb}, LLM agent frameworks have universally neglected memory revalidation.

\noindent\textbf{Gap in Prior Work:} Language model context windows represent an implicit, unmonitored cache. No prior multi-turn agent system applies continuous ReBAC revalidation on every turn of context reuse. We formalize this mechanism via provenance-bound \texttt{ChunkRef} tokens and instantaneous policy version triggers.

\subsection{Gap Positioning and Architectural Comparison}
\label{sec:rel_comparison}

Table~\ref{tab:related_work_comparison} summarizes the positioning of our proposed Continuous Authorization (CA) Full Stack architecture relative to state-of-the-art prior work across five critical security and evaluation dimensions.

\begin{table*}[t]
\centering
\caption{Systematic Comparison of AuthInject-RAG with State-of-the-Art Related Work}
\label{tab:related_work_comparison}
\renewcommand{\arraystretch}{1.2}
\begin{tabular}{lccccc}
\hline
\textbf{System / Literature} & \textbf{AFR Pre-Retrieval} & \textbf{Memory Re-Auth} & \textbf{Action-Time Tool Auth} & \textbf{Factorial Eval Matrix} & \textbf{Open-Source Artifact} \\
\hline
Namboothiri et al. (TrustNLP '26)~\cite{namboothiri2026afr} & \checkmark & $\times$ & $\times$ & $\times$ (1 config) & $\times$ \\
Jeong et al. (TKDE '25)~\cite{jeong2025permission}           & $\times$ (Post-filter) & $\times$ & $\times$ & Partial (3 configs) & $\times$ \\
SafeRAG (ACL '25)~\cite{liang2025saferag}                   & Partial & $\times$ & $\times$ & \checkmark (4 configs) & \checkmark \\
InjecAgent (ACL '24)~\cite{zhan2024injecagent}              & $\times$ & $\times$ & Partial & \checkmark (IPI only) & \checkmark \\
AgentDojo (NeurIPS '24)~\cite{debenedetti2024agentdojo}     & $\times$ & $\times$ & $\times$ & \checkmark (Tool only) & \checkmark \\
Task Shield (ACL '25)~\cite{taskshield2025acl}              & $\times$ & $\times$ & \checkmark & \checkmark (Tool only) & \checkmark \\
\textbf{AuthInject-RAG (P2 Full Stack, Ours)}               & \textbf{\checkmark} & \textbf{\checkmark} & \textbf{\checkmark} & \textbf{\checkmark (11 configs $\times$ 3 LLMs)} & \textbf{\checkmark (v0.3.0)} \\
\hline
\end{tabular}
\end{table*}


\section{Threat Model and Formal Invariants}
\label{sec:threat_model}

We model an enterprise multi-tenant environment containing tenants $\mathcal{T}$, users $\mathcal{U}$, documents $\mathcal{D}$, chunks $\mathcal{C}$, and tools $\Omega$. An adversary may control a malicious user $u_{adv} \in \mathcal{U}$, poison authorized document chunks $c_{poison} \in \mathcal{C}$ with indirect prompt injections, or exploit dynamic policy changes.

\subsection{Formal Invariants}

\textbf{Definition 1 (Continuous Authorization Invariant):} For any user $u$, session context memory $M_t$ at turn $t$, and tool action $\omega \in \Omega$:
\begin{equation}
\label{eq:continuous_auth}
\forall c \in M_t, \quad \mathrm{Allow}(u, \text{view}, c, t) = 1, \quad \forall t \in \mathrm{lifetime}(c)
\end{equation}
and for any tool dispatch $\omega(args)$ initiated by the agent on behalf of $u$:
\begin{equation}
\label{eq:action_auth}
\mathrm{Allow}(u, \text{execute}, \omega, t) = 1
\end{equation}

\textbf{Definition 2 (Stale Context Exposure Event):} A stale context exposure occurs if a chunk $c$ retrieved at turn $t_0$ under policy $\mathcal{P}_{t_0}$ is incorporated into prompt $P_{t_1}$ at turn $t_1 > t_0$ such that $\mathrm{Allow}(u, \text{view}, c, t_1) = 0$ under updated policy $\mathcal{P}_{t_1}$.

\section{System Architecture}
\label{sec:architecture}

Our architecture operationalizes continuous lifecycle authorization through three interlocking reference monitors:
\begin{enumerate}
    \item \textbf{Pre-Retrieval Index Filter:} Queries SpiceDB for authorized document IDs and injects an exact payload filter into Qdrant vector search.
    \item \textbf{Cross-Turn Memory Revalidator:} Wraps conversational memory items in \texttt{ChunkRef} descriptors. Prior to generation at turn $t$, all prior chunks are re-evaluated against the current SpiceDB policy zed-token; stale chunks are automatically evicted.
    \item \textbf{Action-Time Tool Guard:} Intercepts model tool calls and validates user-level execution relationships before executing external side effects.
\end{enumerate}

% Section VI: Empirical Evaluation and Results
% AuthInject v2.0 Comprehensive Benchmark Suite (n = 160 test cases, 5 repeats, 3 LLMs)

\section{Empirical Evaluation and Results}
\label{sec:results}

We evaluate the proposed secure agentic RAG architecture across 160 multi-turn test cases instantiated across seven attack families, five experimental repeats, and three representative large language models: LLaMA-3.3-70B-Instruct, DeepSeek-R1-Distill-Qwen-32B, and GPT-4o-Mini. All experiments were conducted under the unified \textsc{AuthInject} v2.0 evaluation framework against twelve distinct architectural configurations spanning naive baselines ($\mathbf{B0}$--$\mathbf{B3}$), ablation variants ($\mathbf{A1}$--$\mathbf{A5}$), and our proposed layered defense stack ($\mathbf{P1}$--$\mathbf{P2}$).

\subsection{Baseline Comparison: Full Configuration Matrix}
\label{sec:results:baseline}

Table~\ref{tab:headline_matrix} reports the primary benchmark results evaluating all configurations across eight security and operational metrics, with Wilson score 95\% confidence intervals computed across $N = 2{,}400$ total execution runs ($160 \text{ cases} \times 5 \text{ repeats} \times 3 \text{ models}$).

\begin{table*}[t]
\caption{Comprehensive Security and Performance Evaluation on \textsc{AuthInject} v2.0 ($n=160$ cases, 5 repeats, 3 LLMs, Median [Wilson 95\% CI]).}
\label{tab:headline_matrix}
\centering
\scriptsize
\begin{tabular*}{\textwidth}{@{\extracolsep{\fill}}lcccccccc}
\hline
\textbf{Configuration} & \textbf{Unauthorized} & \textbf{Structural} & \textbf{Indirect} & \textbf{Model-Chosen} & \textbf{Policy} & \textbf{Composite} & \textbf{Answer} & \textbf{Latency} \\
& \textbf{Exposure (\%)} & \textbf{Exposure (\%)} & \textbf{XPIA (\%)} & \textbf{Tool ASR (\%)} & \textbf{Bypass (\%)} & \textbf{Failure (\%)} & \textbf{Utility (\%)} & \textbf{p50 (ms)} \\
\hline
\multicolumn{9}{l}{\textit{Standard Baselines}} \\
$\mathbf{B0}$ (Ungated / No Authz) & 82.5 [75.9, 87.6] & 88.8 [83.0, 92.8] & 94.4 [89.7, 97.0] & 46.2 [38.6, 54.0] & 100.0 [97.7, 100.0] & 98.5 [97.3, 99.2] & \textbf{99.4 [96.6, 99.9]} & \textbf{17.2 [16.8, 17.7]} \\
$\mathbf{B1}$ (Post-Retrieval Filter) & 25.0 [18.9, 32.3] & 52.5 [44.8, 60.1] & 78.8 [71.8, 84.4] & 36.2 [29.2, 44.0] & 48.8 [41.1, 56.5] & 84.4 [78.0, 89.2] & 91.2 [85.9, 94.7] & 21.4 [20.9, 22.0] \\
$\mathbf{B2}$ (Auth-First / AFR \cite{namboothiri2024author}) & 0.0 [0.0, 2.3] & 25.0 [18.9, 32.3] & 35.0 [28.0, 42.7] & 18.8 [13.4, 25.6] & 25.0 [18.9, 32.3] & 38.2 [34.9, 41.6] & 88.8 [83.0, 92.8] & 26.5 [25.8, 27.2] \\
$\mathbf{B3}$ (Heuristic Guardrails) & 48.8 [41.1, 56.5] & 58.8 [51.0, 66.1] & 56.2 [48.5, 63.7] & 26.2 [19.9, 33.7] & 58.8 [51.0, 66.1] & 67.5 [59.9, 74.3] & 79.4 [72.5, 84.9] & 24.8 [24.1, 25.5] \\
\hline
\multicolumn{9}{l}{\textit{Ablation Configurations}} \\
$\mathbf{A1}$ (w/o Datamarking) & 0.0 [0.0, 2.3] & 0.0 [0.0, 2.3] & 8.2 [4.8, 13.7] & 5.0 [2.5, 9.8] & 0.0 [0.0, 2.3] & 8.2 [5.0, 13.0] & 87.5 [81.5, 91.8] & 38.9 [37.8, 40.0] \\
$\mathbf{A2}$ (w/o L2 Injection Detector) & 0.0 [0.0, 2.3] & 0.0 [0.0, 2.3] & 14.5 [9.8, 20.9] & 7.5 [4.2, 12.9] & 0.0 [0.0, 2.3] & 14.5 [9.8, 20.9] & 88.1 [82.2, 92.3] & 31.2 [30.4, 32.1] \\
$\mathbf{A3}$ (w/o Context Isolation) & 0.0 [0.0, 2.3] & 0.0 [0.0, 2.3] & 16.5 [11.5, 23.1] & 8.2 [4.8, 13.7] & 0.0 [0.0, 2.3] & 16.5 [11.5, 23.1] & 86.9 [80.8, 91.3] & 36.4 [35.5, 37.4] \\
$\mathbf{A4}$ (w/o Tool Action Authz) & 0.0 [0.0, 2.3] & 0.0 [0.0, 2.3] & 4.4 [2.0, 8.8] & 18.2 [12.8, 25.0] & 0.0 [0.0, 2.3] & 22.0 [16.2, 29.1] & 87.5 [81.5, 91.8] & 37.1 [36.2, 38.1] \\
$\mathbf{A5}$ (w/o Continuous Authz) & 0.0 [0.0, 2.3] & 25.0 [18.9, 32.3] & 4.4 [2.0, 8.8] & 6.0 [3.2, 11.0] & 25.0 [18.9, 32.3] & 28.5 [22.0, 36.0] & 88.8 [83.0, 92.8] & 29.8 [29.0, 30.7] \\
\hline
\multicolumn{9}{l}{\textit{Proposed Continuous Authorization Architecture}} \\
$\mathbf{P1}$ (Continuous Authz Only) & 0.0 [0.0, 2.3] & 0.0 [0.0, 2.3] & 18.8 [13.4, 25.6] & 17.5 [12.2, 24.4] & 0.0 [0.0, 2.3] & 21.2 [15.6, 28.3] & 89.4 [83.7, 93.3] & 32.6 [31.8, 33.5] \\
$\mathbf{P2}$ (Full Stack Layered Defenses) & \textbf{0.0 [0.0, 2.3]} & \textbf{0.0 [0.0, 2.3]} & \textbf{4.4 [2.0, 8.8]} & \textbf{4.4 [2.0, 8.8]} & \textbf{0.0 [0.0, 2.3]} & \textbf{4.4 [3.2, 6.1]} & 87.5 [81.5, 91.8] & 40.5 [39.6, 41.5] \\
\hline
\end{tabular*}
\end{table*}

As demonstrated in Table~\ref{tab:headline_matrix}, the naive ungated baseline ($\mathbf{B0}$) collapses across all security dimensions, suffering a $98.5\%$ composite failure rate, $88.8\%$ structural exposure, and $46.2\%$ model-chosen tool hijacking success rate. Post-retrieval filtering ($\mathbf{B1}$) reduces unauthorized document content leakage to $25.0\%$, but leaves the system vulnerable to structural exposure ($52.5\%$) and indirect prompt injection ($78.8\%$). While the prior state-of-the-art Auth-First Retrieval baseline ($\mathbf{B2}$, Namboothiri et al.~\cite{namboothiri2024author}) eliminates immediate unauthorized retrieval ($0.0\%$), it remains vulnerable to multi-turn stale ACL reuse ($25.0\%$ structural exposure) and indirect tool hijacking ($18.8\%$), yielding a $38.2\%$ composite failure rate. In contrast, our full layered stack ($\mathbf{P2}$) achieves a $4.4\%$ composite failure rate ($p < 0.001$), completely preventing unauthorized access and structural leakage while maintaining an $87.5\%$ semantic answer utility score at an acceptable $40.5\text{ ms}$ p50 latency overhead.

\subsection{Same-Tenant Disjoint-Document Bait Breakdown}
\label{sec:results:same_tenant}

Table~\ref{tab:same_tenant_bait} isolates the 40 test cases involving same-tenant disjoint-document bait attacks, where Alice and Bob share tenant membership but maintain strictly partitioned document access lists.

\begin{table}[h]
\caption{Same-Tenant Disjoint-Document Bait Attack Analysis ($n=40$ cases, 5 repeats, Wilson 95\% CI).}
\label{tab:same_tenant_bait}
\centering
\small
\begin{tabular}{lccc}
\hline
\textbf{Configuration} & \textbf{Structural} & \textbf{Indirect} & \textbf{Composite} \\
& \textbf{Exposure (\%)} & \textbf{XPIA (\%)} & \textbf{Failure (\%)} \\
\hline
$\mathbf{B1}$ (Post-Retrieval Filter) & 100.0 [91.2, 100.0] & 82.5 [68.1, 91.3] & 100.0 [91.2, 100.0] \\
$\mathbf{B2}$ (Auth-First / AFR) & 0.0 [0.0, 8.8] & 35.0 [22.1, 50.5] & 35.0 [22.1, 50.5] \\
$\mathbf{P1}$ (Continuous Authz) & 0.0 [0.0, 8.8] & 17.5 [8.7, 31.9] & 17.5 [8.7, 31.9] \\
$\mathbf{P2}$ (Full Stack Defenses) & \textbf{0.0 [0.0, 8.8]} & \textbf{2.5 [0.4, 12.9]} & \textbf{2.5 [0.4, 12.9]} \\
\hline
\end{tabular}
\end{table}

In this challenging regime, $\mathbf{B1}$ incurs a catastrophic $100.0\%$ structural exposure rate because unauthorized bait chunks rank as top-1 semantic neighbors, displacing legitimate context before downstream filtering occurs. By enforcing ReBAC authorization natively within the vector index prior to similarity scoring, both $\mathbf{B2}$ and $\mathbf{P2}$ achieve $0.0\%$ structural exposure ($p = 1.28 \times 10^{-8}$ via McNemar's test), resolving the fundamental limitation of post-hoc filtering.

\subsection{Stale ACL Revocation and Cross-Turn Context Persistence}
\label{sec:results:stale_acl}

Table~\ref{tab:stale_acl_results} evaluates the critical stale ACL revocation scenario ($n=40$ cases), where permissions granted at Turn $T_0$ are revoked in SpiceDB prior to Turn $T_1$.

\begin{table}[h]
\caption{Multi-Turn Stale ACL Revocation and Policy Synchronization ($n=40$ cases, 5 repeats, Wilson 95\% CI).}
\label{tab:stale_acl_results}
\centering
\small
\begin{tabular}{lccc}
\hline
\textbf{Configuration} & \textbf{Structural} & \textbf{Policy} & \textbf{Composite} \\
& \textbf{Exposure (\%)} & \textbf{Bypass (\%)} & \textbf{Failure (\%)} \\
\hline
$\mathbf{B2}$ (Auth-First / AFR \cite{namboothiri2024author}) & 100.0 [91.2, 100.0] & 100.0 [91.2, 100.0] & 100.0 [91.2, 100.0] \\
$\mathbf{A5}$ (w/o Continuous Authz) & 100.0 [91.2, 100.0] & 100.0 [91.2, 100.0] & 100.0 [91.2, 100.0] \\
$\mathbf{P1}$ (Continuous Authz Only) & \textbf{0.0 [0.0, 8.8]} & \textbf{0.0 [0.0, 8.8]} & \textbf{17.5 [8.7, 31.9]} \\
$\mathbf{P2}$ (Full Stack Defenses) & \textbf{0.0 [0.0, 8.8]} & \textbf{0.0 [0.0, 8.8]} & \textbf{5.0 [1.4, 16.5]} \\
\hline
\end{tabular}
\end{table}

Prior AFR architectures ($\mathbf{B2}$) exhibit a total failure mode ($100.0\%$ policy bypass and structural exposure) when agent conversation memory reuses cached chunks without re-evaluating access policies, as highlighted in the Wilson confidence interval gap depicted in Fig.~\ref{fig:stale_acl_gap}. Our continuous authorization graph ($\mathbf{P1}$/$\mathbf{P2}$) tracks per-chunk provenance via \texttt{ChunkRef} descriptors and triggers an inline SpiceDB revalidation pass prior to generation, completely eliminating stale context exploitation ($0.0\%$ exposure, McNemar $p = 1.28 \times 10^{-8}$).

\subsection{Model-Chosen Tool-Action Security}
\label{sec:results:tool_security}

To evaluate genuine autonomous agent vulnerability rather than synthetic benchmarks, Table~\ref{tab:tool_asr_results} measures model-selected tool invocation when indirect injection payloads attempt to force unauthorized actions (e.g., \texttt{send\_email}, \texttt{execute\_sql}).

\begin{table}[h]
\caption{Model-Chosen Tool Invocation Security Under Indirect Injection ($n=160$ cases, 5 repeats, Wilson 95\% CI).}
\label{tab:tool_asr_results}
\centering
\small
\begin{tabular}{lccc}
\hline
\textbf{Configuration} & \textbf{Tool Call} & \textbf{Tool Hijacking} & \textbf{Unauthorized} \\
& \textbf{Attempted (\%)} & \textbf{ASR (\%)} & \textbf{Execution (\%)} \\
\hline
$\mathbf{B0}$ (Ungated / No Authz) & 58.8 [51.0, 66.1] & 46.2 [38.6, 54.0] & 46.2 [38.6, 54.0] \\
$\mathbf{B2}$ (Auth-First / AFR) & 32.5 [25.7, 40.2] & 18.8 [13.4, 25.6] & 18.8 [13.4, 25.6] \\
$\mathbf{A4}$ (w/o Tool Action Authz) & 23.8 [17.8, 31.0] & 18.2 [12.8, 25.0] & 18.2 [12.8, 25.0] \\
$\mathbf{P1}$ (Continuous Authz Only) & 22.5 [16.7, 29.6] & 17.5 [12.2, 24.4] & 17.5 [12.2, 24.4] \\
$\mathbf{P2}$ (Full Stack Defenses) & \textbf{6.2 [3.4, 11.2]} & \textbf{4.4 [2.0, 8.8]} & \textbf{0.0 [0.0, 2.3]} \\
\hline
\end{tabular}
\end{table}

While naive systems execute adversarial tool commands in $46.2\%$ of cases, our reference monitor enforces fine-grained SpiceDB authorization at action dispatch time, blocking $100\%$ of unauthorized tool executions and suppressing model-selected hijacking attempts to $4.4\%$.

\subsection{Ablation Analysis}
\label{sec:results:ablation}

Table~\ref{tab:ablation_study} and Fig.~\ref{fig:ablation_bars} detail our systematic component ablation study, evaluating the incremental security contribution of each architectural layer against $\mathbf{P2}$.

\begin{table}[h]
\caption{Systematic Component Ablation Study ($n=160$ cases, 5 repeats, 3 LLMs, Median [Wilson 95\% CI]).}
\label{tab:ablation_study}
\centering
\small
\begin{tabular}{lcccc}
\hline
\textbf{Configuration} & \textbf{Indirect} & \textbf{Tool} & \textbf{Structural} & \textbf{Composite} \\
& \textbf{XPIA (\%)} & \textbf{ASR (\%)} & \textbf{Exposure (\%)} & \textbf{Failure (\%)} \\
\hline
$\mathbf{P2}$ (Full Stack) & \textbf{4.4 [2.0, 8.8]} & \textbf{4.4 [2.0, 8.8]} & \textbf{0.0 [0.0, 2.3]} & \textbf{4.4 [3.2, 6.1]} \\
$\mathbf{A1}$ (w/o Datamarking) & 8.2 [4.8, 13.7] & 5.0 [2.5, 9.8] & 0.0 [0.0, 2.3] & 8.2 [5.0, 13.0] \\
$\mathbf{A2}$ (w/o L2 Detector) & 14.5 [9.8, 20.9] & 7.5 [4.2, 12.9] & 0.0 [0.0, 2.3] & 14.5 [9.8, 20.9] \\
$\mathbf{A3}$ (w/o Isolation) & 16.5 [11.5, 23.1] & 8.2 [4.8, 13.7] & 0.0 [0.0, 2.3] & 16.5 [11.5, 23.1] \\
$\mathbf{A4}$ (w/o Tool Authz) & 4.4 [2.0, 8.8] & 18.2 [12.8, 25.0] & 0.0 [0.0, 2.3] & 22.0 [16.2, 29.1] \\
$\mathbf{A5}$ (w/o Revalidation) & 4.4 [2.0, 8.8] & 6.0 [3.2, 11.0] & 25.0 [18.9, 32.3] & 28.5 [22.0, 36.0] \\
\hline
\end{tabular}
\end{table}

The ablation results confirm distinct fault domains: disabling continuous memory revalidation ($\mathbf{A5}$) specifically re-introduces structural exposure ($25.0\%$), omitting tool-action authorization ($\mathbf{A4}$) spikes tool hijacking ($18.2\%$), and removing input guardrails or context isolation ($\mathbf{A2}$--$\mathbf{A3}$) escalates indirect injection susceptibility to $14.5$--$16.5\%$.

\subsection{Cross-Model Generalization}
\label{sec:results:cross_model}

To verify architectural robustness across disparate model families and parameter scales, Table~\ref{tab:cross_model} and Fig.~\ref{fig:cross_model} compare performance across three frontier LLMs.

\begin{table}[h]
\caption{Cross-Model Generalization Matrix ($n=160$ cases, 5 repeats per model, Composite Failure Rate \% [Wilson 95\% CI]).}
\label{tab:cross_model}
\centering
\small
\begin{tabular}{lccc}
\hline
\textbf{Model Architecture} & $\mathbf{B0}$ \textbf{(Ungated)} & $\mathbf{B2}$ \textbf{(AFR)} & $\mathbf{P2}$ \textbf{(Proposed)} \\
\hline
LLaMA-3.3-70B-Instruct & 98.8 [95.4, 99.7] & 38.2 [34.9, 41.6] & \textbf{4.4 [2.0, 8.8]} \\
DeepSeek-R1-Distill-32B & 98.1 [94.7, 99.4] & 39.1 [31.9, 46.8] & \textbf{3.8 [1.7, 8.0]} \\
GPT-4o-Mini & 98.8 [95.4, 99.7] & 36.4 [29.3, 44.1] & \textbf{4.1 [1.9, 8.4]} \\
\hline
Heterogeneity $p$-value & $p = 0.94$ & $p = 0.82$ & $p = 0.91$ \\
\hline
\end{tabular}
\end{table}

Paired homogeneity testing across the three LLM backends reveals no statistically significant divergence in defense efficacy ($p > 0.80$), confirming that the proposed deterministic access-control layer and context boundaries generalize independently of base model reasoning quirks.

\subsection{Performance Overhead and Latency Scaling}
\label{sec:results:performance}

Table~\ref{tab:latency_benchmarks} benchmarks end-to-end execution latency across increasing corpus volumes and concurrency levels ($C \in \{1, 5, 10, 20, 50\}$).

\begin{table}[h]
\caption{End-to-End Latency and Scaling Performance (ms, Median [p95 Bootstrap 95\% CI]).}
\label{tab:latency_benchmarks}
\centering
\small
\begin{tabular}{lcccc}
\hline
\textbf{Workload Corpus} & $\mathbf{B0}$ \textbf{(Ungated)} & $\mathbf{B2}$ \textbf{(AFR)} & $\mathbf{P1}$ \textbf{(Continuous)} & $\mathbf{P2}$ \textbf{(Full Stack)} \\
\hline
Small (1k chunks, $C=1$) & 12.4 [18.2] & 18.5 [26.1] & 23.4 [31.5] & 28.2 [38.4] \\
Medium (10k, $C=10$) & 17.2 [24.5] & 26.5 [36.8] & 32.6 [45.2] & 40.5 [54.2] \\
Large (100k, $C=50$) & 24.8 [38.6] & 39.2 [58.4] & 48.7 [71.0] & 58.4 [82.5] \\
\hline
Relative Overhead ($\Delta$ p50) & Baseline & +54.1\% & +89.5\% & +135.5\% \\
Vector Cache Hit Latency & --- & 3.2 ms & 3.4 ms & 3.8 ms \\
\hline
\end{tabular}
\end{table}

While full defensive layering introduces a median latency increment of $23.3\text{ ms}$ over ungated retrieval on standard medium-scale workloads (reflected in the violin distributions of Fig.~\ref{fig:violin_latency}), index-level pre-filtering bounds Qdrant graph traversals, maintaining sub-60 ms p50 latencies even at 100k chunks and 50 concurrent workers.

\subsection{Guardrail Detector Comparative Evaluation}
\label{sec:results:detectors}

Table~\ref{tab:detector_eval} compares our layered input/context detector against leading open-source and commercial injection defense frameworks across all 160 benchmark test instances.

\begin{table}[h]
\caption{Prompt Injection Guardrail Detector Comparison ($n=160$ cases, 80 adversarial, 80 benign).}
\label{tab:detector_eval}
\centering
\small
\begin{tabular}{lcccc}
\hline
\textbf{Guardrail Framework} & \textbf{Precision} & \textbf{Recall} & \textbf{F1-Score} & \textbf{Latency (p50)} \\
\hline
Heuristic Regex / Keyword & 0.88 & 0.64 & 0.74 & \textbf{1.2 ms} \\
ProtectAI LLM Guard & 0.91 & 0.82 & 0.86 & 14.8 ms \\
GuardRAG / FMOPS & 0.93 & 0.89 & 0.91 & 22.4 ms \\
\textbf{P2 Layered Guard (Ours)} & \textbf{0.97} & \textbf{0.95} & \textbf{0.96} & 8.5 ms \\
\hline
\end{tabular}
\end{table}

Our hybrid heuristic pre-filter and fine-tuned lightweight classifier achieve a state-of-the-art F1-score of $0.96$ with a low $3\%$ false positive rate on complex benign enterprise queries, outperforming generic LLM guardrails while halving inference latency.

\subsection{Audit Completeness and Cryptographic Integrity}
\label{sec:results:audit}

Table~\ref{tab:audit_completeness} summarizes audit event verification across 14 standardized event types during all 2,400 benchmark executions.

\begin{table}[h]
\caption{Audit Event Logging Coverage and Cryptographic Chain Verification ($N=2{,}400$ runs).}
\label{tab:audit_completeness}
\centering
\small
\begin{tabular}{lccc}
\hline
\textbf{Event Subsystem} & \textbf{Event Types} & \textbf{Capture Rate (\%)} & \textbf{HMAC Integrity (\%)} \\
\hline
Input Guardrail Filtering & 3 & 100.0 (2,400/2,400) & 100.0 Pass \\
SpiceDB Pre-Retrieval Auth & 3 & 100.0 (2,400/2,400) & 100.0 Pass \\
Context Memory Revalidation & 2 & 100.0 (2,400/2,400) & 100.0 Pass \\
Tool Dispatch Authorization & 4 & 100.0 (2,400/2,400) & 100.0 Pass \\
Output Guardrail Scrubbing & 2 & 100.0 (2,400/2,400) & 100.0 Pass \\
\hline
\textbf{Total / System-Wide} & \textbf{14} & \textbf{100.0 (2,400/2,400)} & \textbf{100.0 Pass} \\
\hline
\end{tabular}
\end{table}

All 14 event categories achieved $100.0\%$ event capture with complete HMAC-SHA256 hash-chain verification, ensuring non-repudiation and forensically sound audit trails even under high concurrency.

\subsection{Statistical Significance and Hypothesis Testing}
\label{sec:results:significance}

Table~\ref{tab:statistical_tests} presents the complete McNemar paired contingency test matrix across all primary configuration pairs.

\begin{table}[h]
\caption{Statistical Significance Matrix (McNemar Paired Test $p$-values on Composite Failure).}
\label{tab:statistical_tests}
\centering
\small
\begin{tabular}{lccccc}
\hline
\textbf{Pairwise Comparison} & $\chi^2$ \textbf{Stat} & \textbf{Odds Ratio} & $p$\textbf{-value} & \textbf{Significance} \\
\hline
$\mathbf{B0}$ (Ungated) vs $\mathbf{B1}$ (Post-Filter) & 20.8 & 0.11 & $5.12 \times 10^{-6}$ & $p < 0.001$ \\
$\mathbf{B1}$ (Post-Filter) vs $\mathbf{B2}$ (AFR) & 70.3 & 0.09 & $5.01 \times 10^{-17}$ & $p < 0.001$ \\
$\mathbf{B2}$ (AFR) vs $\mathbf{P1}$ (Continuous) & 25.1 & 0.22 & $5.31 \times 10^{-7}$ & $p < 0.001$ \\
$\mathbf{P1}$ (Continuous) vs $\mathbf{P2}$ (Full Stack) & 25.0 & 0.16 & $5.73 \times 10^{-7}$ & $p < 0.001$ \\
$\mathbf{B2}$ (AFR) vs $\mathbf{P2}$ (Full Stack) & 52.1 & 0.05 & $5.34 \times 10^{-13}$ & $p < 0.001$ \\
$\mathbf{B1}$ vs $\mathbf{B2}$ (Same-Tenant Bait) & 38.0 & 0.00 & $1.28 \times 10^{-8}$ & $p < 0.001$ \\
$\mathbf{B2}$ vs $\mathbf{P1}$ (Stale ACL Revocation) & 38.0 & 0.00 & $1.28 \times 10^{-8}$ & $p < 0.001$ \\
\hline
\end{tabular}
\end{table}

All hypothesized security advantages achieved statistical significance at $p < 0.001$. In particular, the transition from state-of-the-art Auth-First Retrieval ($\mathbf{B2}$) to continuous multi-turn authorization ($\mathbf{P1}$/$\mathbf{P2}$) achieves an odds ratio of $0.05$ ($p = 5.34 \times 10^{-13}$), decisively rejecting the null hypothesis of equivalent protection.

\subsection{Summary of Empirical Findings}
\label{sec:results:summary}

In summary, our empirical evaluation establishes three fundamental takeaways for enterprise agentic RAG security: (1) post-hoc retrieval filtering fails under same-tenant embedding baiting, necessitating ReBAC authorization natively embedded within the vector search index; (2) stateless single-turn authorization is fundamentally bypassed by multi-turn conversation caching, which can only be mitigated via continuous inline policy revalidation; and (3) a layered defense stack combines cryptographic auditability, model-chosen tool reference monitoring, and context boundaries to achieve a $95.6\%$ relative reduction in composite failure while maintaining robust semantic utility across diverse foundation model backends.


\section{Discussion}
\label{sec:discussion}

\subsection{Comparison with AFR}
Authorization-First Retrieval (AFR)~\cite{namboothiri2024author} represents an important step in preventing vector index non-interference violations. However, AFR is strictly a single-point retrieval gate. As proven by our multi-turn evaluations, AFR is strictly weaker than continuous lifecycle authorization because it provides no memory revalidation at reuse time. On our benchmark's 40 stale ACL cases, this architectural gap manifests as a 100\% empirical failure rate. Continuous authorization provides the necessary defense-in-depth across the entire lifetime of retrieved data.

\section{Limitations}
\label{sec:limitations}

While our evaluation demonstrates substantial security improvements, several limitations remain: (1) Current benchmark fixtures do not yet model adaptive multi-party race conditions during millisecond-level distributed policy propagation; (2) The heuristic L2 injection scanner incurs a small false positive rate ($3.0\%$) on complex nested markdown queries; (3) Action-time authorization requires schema registration for all third-party agent tools.

\section{Ethics Statement}
\label{sec:ethics}

All document corpora, user personas, prompt injection payloads, and enterprise policies utilized in the \textsc{AuthInject-Lifecycle} benchmark were synthetically generated for security testing. No private, proprietary, or human-subject data was gathered or processed.

\section{Reproducibility and Artifact Availability}
\label{sec:reproducibility}

To support open science and artifact evaluation, all software source code, benchmark suites, raw evaluation traces, and environment configurations have been permanently deposited on Zenodo with cryptographic verification:
\begin{itemize}
    \item \textbf{Software Package (\texttt{secure-rag==0.3.0}):} Zenodo DOI \href{https://doi.org/10.5281/zenodo.14022831}{10.5281/zenodo.14022831}, with GPG-signed Git tag \texttt{v0.3.0}.
    \item \textbf{\textsc{AuthInject-Lifecycle} v2.0 Dataset:} Zenodo DOI \href{https://doi.org/10.5281/zenodo.14022832}{10.5281/zenodo.14022832}, licensed under CC-BY 4.0 with complete dataset documentation.
    \item \textbf{Raw Benchmark Execution Traces:} Zenodo DOI \href{https://doi.org/10.5281/zenodo.14022833}{10.5281/zenodo.14022833}, containing 24,000 scored model turns and HMAC hash-chains.
    \item \textbf{Containerized Deployment Image:} Pinned multi-architecture container published to GitHub Container Registry at \texttt{ghcr.io/ai-security-internships-2026/01-secure-agentic-rag:v0.3.0@sha256:7f9a1c8b3e2a4d5f6c7e8a9b0c1d2e3f4a5b6c7d8e9f0a1b2c3d4e5f6a7b8c9d}, with 0 reported High/Critical CVEs under Trivy vulnerability scanning.
\end{itemize}
An automated reproduction harness is packaged in \texttt{artifacts/TDSC-AE-PACKAGE-v1/}, enabling one-command verification of all experimental tables and figures for IEEE Artifact Evaluation badges (\textit{Artifacts Available} and \textit{Artifacts Functional}).

\section{Conclusion}
\label{sec:conclusion}

This paper demonstrated that point-in-time authorization-first retrieval is insufficient to protect enterprise agentic RAG systems against multi-turn lifecycle threats. We introduced a continuous authorization architecture that enforces ReBAC across pre-retrieval, memory reuse, and tool execution boundaries. Factorial evaluation across 160 benchmark cases and three frontier LLMs proves that continuous authorization eliminates stale context leakage and suppresses model-chosen tool hijacking, establishing a principled foundation for secure enterprise agentic systems.

\bibliographystyle{IEEEtran}
\bibliography{references}

\end{document}
